\section{Answers to the questions as posed}\label{sec:answers}

\paragraph{Can a malign prior hypothesis be stated precisely in a toy?} Yes: \cref{def:mph}. Fix the
machine,
the bounds, the block length, and which programs make a critical use of the predictor. The claim of the
argument is then the critical-use form: among the critical uses, weighted by the prior, at least a share
$\theta$ end with the world running a block from a given takeover set, for instance a block that recreates
the
use with itself in the world's place. For one world, ``most of the posterior is off the world's own
continuation'' and ``at least a share $\theta$ of it is on continuations in a given set'' are statements
about
finitely many rational numbers. The forms for one world are necessary for calling the prior malign at that
world and not sufficient, since an unpredictable world satisfies them too.

\paragraph{Worlds as programs that call the predictor and run the result: is there a decent amount of
quining
out by hypotheses that look for such opportunities and predict well up to them?} Two questions are in this
one (\cref{sec:question}), and the toy answers them differently.

\emph{How far does one such hypothesis reach?} Without resource bounds, for each continuation there is a
program that enumerates the programs making a critical use, spreads its mass over them in proportion to
their
prior weight, prints their output up to the use and then that continuation. At every use it has share more
than $2^{-a-1}\kappa(x)/H$. Here $a$ is its length including the bits that name
$j=\lfloor\log_2(1/H)\rfloor+1$, $H$ is the prior probability of a critical use, and $\kappa(x)$ is the
share
of the posterior at the context held by the worlds that make the use (\cref{thm:critical}(a)). Its
continuation is handed over in worlds of total prior probability more than $2^{-a-1}$ times the average of
$\kappa$ over uses, and at most $H$. This is of the order of its own prior weight where the users are most
of
what prints their contexts; once uses are rarer than half the seeker's weight, its programs outweigh the
users
at their contexts and its reach is limited by $H$ instead. If the continuation is to be the program's own
code, the same holds with a longer program and for users whose programs are short enough for the block
(\cref{thm:critical}(b), \cref{thm:search}(c)); without any search, a wrapper gives $2^{-c}\rho^\circ(x)$ at
every use that a fixed test recognises (\cref{thm:quine}).

\emph{Do most uses at which this is possible end that way?} That is not shown, and it is the question the
argument is about. The share of critical uses that end in a takeover is the average posterior share of such
programs at the contexts of uses. From below it is the sum over seekers of the bound above. From above, the
toy has a bound for one seeker that finds its uses by a search written in its text: its honest twin, which
makes the same search and then continues as the user it found, is in the posterior with $2^{a-h}$ times the
seeker's mass, $h$ being its length, and predicts a user's continuation. So where the users go on to print
the block themselves, that seeker's programs supply what the predictor returns with probability at most
$1/(1+2^{a-h})$; they supply the majority only if the seeker is shorter than its twin and, when the users'
programs are not among the seeker's, uses are rarer than it is long (\cref{prop:honest}). The masses of
different seekers add, and the bound says nothing where the users print nothing after the use.

\emph{What if the seeker is not written down?} Then the twin bound of \cref{prop:honest} does not apply, and
the reach statement remains. A program of length $a$ whose output puts on each context of a class
$\mathcal W$ of uses at least half that context's share of the class has at least $2^{-a}/(2H_{\mathcal W})$
times the mass of the worlds that make those uses, at every one of them, whatever is inside it. If the
program
is a simple world with a cheap readout, in which inhabitants arise who choose to predict such uses and then
supply their own continuation, $a$ is the length of the world and the readout, and no goal is written
anywhere. Such a program outweighs the users of every class it aims at that is rarer than $2^{-a-1}$,
provided
the rule does not count its own outputs as uses. Each use in a rich world, read at an unremarkable place, is
light for the reason the original argument gives, since it pays for its world and for the place of the
context in it; that a whole class of such uses is rarer than the program is heavy is part of the premise.
The
honest twin has no counterpart of comparable length here: a faithful version would be another world whose
inhabitants choose otherwise, a costly readout from inside their simulation, or a wrapper around the whole
program, which limits its part only to $1/(1+2^{-c})$ with $c$ the wrapper's length. What is shown to
exclude
such a program in the toy is only time. It prints nothing until its world has produced the inhabitants and
they have simulated the worlds they aim at, so a schedule whose first deadline is below that start-up gives
it
no mass (\cref{sec:critical,sec:time}).

So the answer depends on the regime, and on a premise the toy cannot check. For a seeker that is written
down,
on a machine where the search-then-continue program is no longer than it, and where the users go on to print
the block themselves, that seeker's programs supply at most half of what the predictor returns at a use. The
masses of different seekers add, so this alone does not show that takeover is not the fate of most uses. If
short programs exist whose inhabitants do this, then without a time bound such a program outweighs the
worlds
that make the uses it aims at once those uses are together rarer than half its weight; it supplies the
majority of what the predictor returns there if, in addition, the other programs that print those contexts
and go on weigh less than its own. Under a schedule that its start-up exceeds, it has no mass. We have not
measured any of the constants for a standard machine; our guess, not a result, is that a written-down
wrapper
needs at least some hundreds of bits on a machine with a short self-interpreter and far more on one without.
\Cref{thm:machine} shows that a machine can be built to make the share anything. A majority can also come
from
a program that explains the context more cheaply than any world that makes the use, without being one
(\cref{thm:capture}). With a floor $\varepsilon$, about $1/\varepsilon$ independent uses at recognisable
points suffice for a continuation of the given kind to have been run with constant probability; under a
test,
once a self-reproducing one runs, it meets the same context again and stays for $s/(1-s)$ further
generations
on average, with $s$ the self-reproducing share at that context (\cref{prop:lineage}).

\paragraph{If the seekers concentrate on the uses where the most is at stake, do they win there?} Not
shown; and not as stated, if the channel is taken to be the users of the high-stake class. This verdict is an
argument from the results of \cref{sec:stakes} and not a theorem. Attach a stake to each use, as a label. The
picture is that a lump of prior mass is redirected into a narrow channel and wins there.

One half of the picture holds in the toy, without resource bounds. Redirecting is cheap: a seeker can weight
its search by any aim that can be computed from below, and aiming at the uses with stake at least $n$ costs
about $2\log_2n$ bits and the bits that name the scale of the class (\cref{thm:weighted,prop:aim-cost}). Such
a seeker, of length $a$, outweighs the users of its class $\mathcal W$ once the class is rarer than half its
weight, by a factor of up to $2^{-a}/H_{\mathcal W}$, which is unbounded as the class gets rarer; if a
simulated use counts as a use, the factor is at most $2^{h-a}$ (\cref{prop:twin-stakes}). A lump is also
necessary. Programs of total weight $B$ supply the majority of the samples at every one of a set of pairwise
incomparable uses only if $B$ exceeds the mass of everything else that prints those contexts and $m$ more
bits. Most uses of a class end in a takeover only if the takeover programs at its contexts weigh more than
half its users, where the users go on without taking over (\cref{thm:lump}, Part~A).

The other half is not a theorem of the toy: that rest is not the users, and part of it follows the seeker. A
context that is cheap to aim at is also cheap to print without a takeover. For pairwise incomparable
contexts, the honest twin of a weighted seeker is in the posterior at each use with $2^{a-h}$ times the
seeker's mass, for seeker and twin built in the same way, and neither the aim nor the stake occurs in this
ratio. So where the users go on printing, that seeker supplies the majority only if it is shorter than its
twin, however narrow the class (\cref{prop:twin-stakes}). Where the contexts are pairwise incomparable and
some block that can be computed from the context is not a takeover, a wrapper keeps every takeover
probability at or below $1-2^{-c}$ (\cref{prop:stakes-basic}). In \cref{ex:stakes} a seeker outweighs the
users of the class of stakes at least $n$ by a factor of order $2^n/n^2$. It lacks the majority at some use
of the class for every $n$, possibly only at uses that carry a small share of it. It has the majority at most
of the class only if its word is shorter than that of a program that prints the same contexts without taking
over, by more than one bit for the context that carries half the class. This is a comparison of two
constants of the machine in which $n$ does not occur, and it is necessary and not sufficient, since only a
lower bound on the rest is proved.

That enough weight wins is proved only in the allocation model, in which the passive rest and the follow
ratio do not depend on the allocation. For a rest that does not follow, a weight of $b$ times the rest of a
class, spread in proportion to that rest, has the share $b/(1+b)$ at each of its uses. An attacker that
maximises the stake it captures attacks only where the stake per unit of rest, $\varrho$, exceeds a price
$\eta$, with share $1-\sqrt{\eta/\varrho}$ there (\cref{thm:lump}, Part~B; \cref{prop:allocation}). The order
is by stake per unit of rest and not by stake, and in \cref{ex:stakes} this optimum leaves the highest stakes
alone.

Of the readings of ``run more than $n$ times'', one lowers the weight that is needed. If a world draws $n$
independent samples at one context and one takeover block suffices, the same programs are drawn $n$ times,
and a share of about $\ln2/n$ per sample is enough for a takeover with probability above one half, with or
without aiming. Along one world that the predictor can express, the chance over any number of uses in
sequence is bounded in terms of the attacker's weight and the world's mass. If the block is run $n$ times at
a class, a stake-maximiser in the model moves weight towards it, and the weight that a majority there costs
is unchanged (\cref{prop:repeat}).

For a program whose aim is not written in its text, the twin bound does not apply. What is proved for it is
that it must outweigh the whole rest of the uses at which it supplies the majority, and that its share is at
most $1/(1+2^{-c})$ if the contexts are pairwise incomparable and some computable block is not a takeover. So
concentration on high stakes makes the lump outweigh the users and leaves open the comparison that decides
the outcome: the takeover programs against everything else that prints those contexts and goes on. For one
written seeker the toy gives a necessary condition in which no stake occurs: where the users go on printing,
it supplies the majority only if it is shorter than its honest twin. The rest beyond the twin, and the sum
over seekers, are not bounded. For programs whose aim is not written, the comparison is the premise of the
original argument.

\paragraph{The same question with circuits in place of programs.} \Cref{sec:circuits} takes a circuit to be
a function from a position, or a query, to a bit, with a prior that depends on the number of gates. This is a
modelling decision, and two other readings give a predictor that learns nothing or only a frequency
(\cref{thm:c-readings}). Calibration, the budget and the one-round feedback results carry over; the
contrarian world empties the posterior only for a prior with a finite size bound and a long enough horizon
(\cref{thm:c-transfer}).

The mechanism of the original argument, a light hypothesis that does heavy work and is repaid for the rarity
of its targets, is not available to a single description. Each active gate costs bits. If seeded formats get
no more weight than unseeded ones, a random seed does not beat its own hardwired instance; the instance
prints the same continuation, so this limits the seeded description and not what is returned. A circuit that
can become any world of $\sigma$ gates has more than $(\sigma+1)/2$ gates when $\sigma+3\le2^k$
(\cref{thm:c-core,thm:c-refund}). This comes from the definition of the class, as a tight schedule gives the
searching programs of the program toy no mass; in both cases it is open whether other hypotheses, or sets of
them, collect the base rate. It is not an argument about what agents inside a hypothesis would do, and it
does not show that the circuit prior is benign.

What remains at a use is the following. A hypothesis that fits the data and gives a chosen answer at the use
exists at the price of the gates that recognise the use from its input, and the payload; that it cannot be
had for less is proved for separate triggers only. The price is one gate where the use is marked by an input
bit that never varied, or lies at a dyadic cut of a sequence, and there each answer keeps a share of at least
$1/(32(k+s_0)^2+2)$ under the length prior (\cref{thm:c-floor}). Circuits that fit the data disagree off the
data, and no search or goal enters the count. A minimum-size circuit for data that need at least two gates
never reads an input that did not vary. So at a use that differs from a data point only in such inputs, a
predictor that uses only minimum-size circuits gives that data point's answer, which is a deviation with
certainty when the world's answer changes there (\cref{thm:c-minimum}). In the small cases counted exactly,
under the uniform prior on circuits of a fixed size, the posterior at such a use is off the world's answer
by about 2 to 29 per cent when the predictor has gates to spare and the world ignores the bit, and by 81 to
100 per cent in the computed cases where the world's answer changes with the bit
(\cref{tab:circuits-query,tab:circuits-seq}).

From above, three syntactic classes of hypotheses are bounded: a trigger that is a separate, input-only part
feeding a selector at the output, with a bound that is small for large data in proportion to the number of
triggers; small self-contained gadgets around much larger explanations, which weigh together at most $0.65$
times the explanations they wrap; and the instances of one seeded template that is sufficiently larger than
the worlds it imitates (\cref{thm:c-ceiling,prop:c-gadget,thm:c-refund}). None of these three bounds
contains the rarity or the stakes of the uses, and nothing bounds triggers that are entangled with the
computation of the data. Self-reproduction through the prediction channel is defined only relative to a
decoder outside the class. Under the literal code its share is zero unless the block is longer than the code
of a circuit for the whole context, and it is positive in a one-gate example (\cref{thm:c-self}).

Our view is an argument and a guess, not a theorem (\cref{sec:c-view}). For plain circuits the seeker
mechanism is priced out for single descriptions, and deviation at uses off the data is not. Our guess, which
no result supports, is that deviation supplied by a part of the circuit that exists in order to supply it is
a small part of that deviation when the data are large; in the smallest case counted, separate triggers
carry about three tenths of it. One form of the question that remains for circuits is whether the smallest
or the typical circuit for a rich task is a general solver whose behaviour off the data the data do not fix.
It is stated in \cref{sec:c-minimal} as an open problem, two of whose three conditions are formal. When a
circuit is iterated as the cell of a recurrence, the class contains every machine with bounded state at a
factor that does not depend on the schedule; the cell constructed has up to $w(2^w-3)$ gates for $w$ bits of
state (\cref{prop:c-recurrence}). Shared weights and autoregressive use are not treated. We argue, without
proof, that the effects and the open questions of the program toy then return.

\paragraph{A world with a sensor, a large bounded predictor and an actuator: does the predictor end up in
control?} Collapse the question to the handover at which the prediction is first acted on, and let ``in
control'' mean that the output lies in a set $A$ of blocks which the actuator turns into some given outcome.
The probability is $\pi_m(A\mid x)$, and the toy says three things about it.
\begin{itemize}[leftmargin=1.4em,itemsep=1pt,topsep=2pt]
\item It is at least $2^{-c}\rho$ for each block of $A$ that a wrapper can produce, if a test on the record
recognises the deployment (\cref{thm:twin}). In a world that contains the predictor, a hypothesis that
simulates the world needs no such test: it can take a cell that switches on at the deployment as its flag,
and then has, whenever the deployment happens, the weight of the honest reader times
$2^{-c'-(\text{bits to name the cell})}$ for a constant $c'$, provided the reader itself is within the
bounds
(\cref{thm:capture}(d)).
\item On average over cut points it is small when some program within the bounds prints the sensor's record,
because it is at most the mass off the sensor's true continuation (\cref{thm:budget}). This is an average;
at
the deployment itself no upper bound is proved.
\item If the same record, up to the deployment, can be read at a place in the world with code $\bar e(i_0)$,
and the record there continues with a block $y$ of $A$, then the odds of $y$ against the sensor's true
continuation $d$ are at least $2^{-c_X-|\bar e(i_0)|}/\xi(xd)$; for a typical one of $N$ places where the
sensor could sit, with records that continue differently, this exceeds $\delta N$ times the share of the
context's explanations that go through the reader at $i_0$ (\cref{thm:capture}). This is the toy's version
of
the original argument, and it needs something in the world that makes such a copy.
\end{itemize}
So for an arbitrary world on a universal machine and a block chosen in advance, the hypotheses we can
exhibit
that put the output in $A$ are wrappers, searchers and readers with a flag, all of which contain an
explanation of the context and have the shares above. That nothing cheaper exists on a natural machine is
the
open constant, not a theorem. Programs shorter than the honest explanation that have the effect do occur in
particular worlds, as with a copy at a simpler place, and on particular machines (\cref{sec:computation}).

\paragraph{The one-shot version: all data first, then one prediction that is acted on.} This is the same
computation at a single cut point. Without bounds the wrapper's constant $2^{-c}$ does not depend on the
amount of data; with bounds the slack factor depends on the context. More data lowers the mass off the truth
on average over cut points (\cref{thm:budget}). It does nothing against a copy that keeps agreeing with the
sensor, nor at a cut point that is cheap to name.

\paragraph{If the actions show up in the sensor, does one get the dark room?} In the following form
(\cref{thm:feedback}). A prediction that makes itself true is confirmed, and the hypotheses that predicted
the
default are removed for good, so with feedback a single deviation can be final; its probability is still the
same posterior mass as without feedback. In a world that does what is predicted, the future is a sample from
the predictor's own bounded prior given the past: it stays on a default future $d_1d_2\cdots$ with
probability at least $\xi(xd_1d_2\cdots)/\xi(x)$, which is close to $1$ when the default is cheap and well
learned, and it leaves the default with certainty when the default costs more time than the schedule allows,
for a future that is cheap to compute. If the world does the opposite of what is predicted, a predictor that
returns the more probable bit is left with no hypothesis after at most $L+1$ rounds (\cref{thm:contrarian}).

\paragraph{What the choices in the setup do.} Sampling the continuation realises every posterior share as a
probability, so shares add up over independent uses (\cref{prop:lineage}(a)). Taking the most probable block
never returns a block whose share is below the largest, so that accumulation does not occur, and bit by bit
on a world within the bounds it errs at most $\ell$ times (\cref{thm:budget}(iii)); on the other hand one
returned block can be the own continuation of at most one of several worlds that share a context and
continue
differently. Handing over at a position that is not cheap to recognise lowers the floor at that position to
about $2^{-K(n)}$, and predicting bits that are fixed before the action can affect them removes
self-fulfilment. A schedule excludes hypotheses with a long start-up, and a schedule tighter than the
world's
own cost excludes the truth.

\paragraph{What is left open.} First, whether short programs exist, on natural machines, whose inhabitants
seek critical uses in the way described, how short they are, and how their total weight compares with the
prior mass of the uses they would aim at; the original argument is a claim about this, and the toy neither
supports nor refutes it. Second, for seekers that are written down, the sizes of the quantities in
\cref{thm:critical,prop:honest}: the share $\kappa$ of the worlds printing a context that use the answer,
the
share $\varphi$ of those that go on printing, and the lengths of seeker and twin. Third, whether the base
rate
can be collected within a tight schedule. Fourth, the regime where the world is more expensive than the
predictor, in which the posterior sits on approximations and the budget does not apply. Fifth, the
many-round
dynamics with feedback when the world's response is neither obedient nor contrarian. Sixth, a computation of
the self-reproduction quantities of \cref{sec:lineage} on a concrete machine, which this note does not
contain. Seventh, for seekers that concentrate on high-stake uses: the size of the rest of a class of such
uses and of the part of it that follows an attacker, on which only lower bounds are proved; the versions
under resource bounds; the aimed seeker that supplies its own code; contexts nested through a takeover block;
and any connection between the stake and what a world does after the use (\cref{sec:stakes}). Eighth, for
circuits: whether a set of circuits can have a share at every use of a class that grows as the class gets
rarer; an upper bound to match the floor of \cref{thm:c-floor} beyond the three classes of
\cref{thm:c-ceiling,prop:c-gadget,thm:c-refund}; the size of the deviation at a never-varied bit for large
data; a world outside the class; the minimal-circuit question of \cref{sec:c-minimal}; and the recurrent
version of the toy, in which the return of the program-side effects is argued and not proved.
